\section{Consideraciones de escalabilidad y rendimiento}

La muestra sirve para validar el diseño, no para medir rendimiento. La tabla más
grande tiene 56 filas y \texttt{embedding} tiene 36. Los tiempos y planes
observados no se pueden extrapolar a un sistema real.

Las tablas que más crecerían son \texttt{evento\_auditoria}, \texttt{consulta},
\texttt{respuesta}, \texttt{fragmento} y \texttt{embedding}. La auditoría suma
una fila por cada acción y no se elimina. Por eso es la primera candidata para
un particionado mensual o trimestral, según el volumen y la retención necesaria.

Los fragmentos y embeddings crecen con cada documento o versión. Un fragmento
puede tener un embedding por modelo, lo que permite conservar vectores
históricos. Durante una migración de modelo, el volumen aumenta de forma
temporal. Una política de retención puede limitar ese crecimiento.

Las consultas críticas son las búsquedas top-k y las híbridas autorizadas. Estas
últimas combinan distancia coseno, estado, vigencia y permisos. Se creó un índice
HNSW con \path{vector_cosine_ops} para el crecimiento previsto. En la
muestra, el optimizador elige una exploración secuencial (\texttt{Seq Scan})
porque recorrer 36 vectores es más barato. El plan está documentado en
\path{evidencias/planes_ejecucion.md}.

Los índices B-tree cubren los filtros y relaciones de las consultas: condiciones
comerciales, entregas, incidencias, pedidos, vigencia y auditoría. No se indexa
cada columna, porque los índices ocupan espacio y encarecen las escrituras.

Primero se aumentarían los recursos de PostgreSQL y se revisarían los planes de
ejecución. Las réplicas de lectura sólo se justifican si crece el tráfico y las
mediciones muestran un límite. La publicación, los permisos y la auditoría deben
quedar en la instancia principal para conservar la consistencia transaccional.

Más adelante podrían separarse algunos componentes: archivos en almacenamiento
de objetos, auditoría particionada o un motor orientado a escritura masiva. Una
base vectorial dedicada sólo se justifica si los embeddings y el volumen de
búsquedas por segundo crecen varios órdenes de magnitud. Esa opción obliga a
sincronizar permisos y vigencia entre motores, con más complejidad y riesgo de
inconsistencia.

La decisión debe basarse en crecimiento por tabla, latencia del top-k, costo de
auditoría y tolerancia a consistencia diferida. Hasta que esas mediciones muestren
un límite, PostgreSQL conserva las garantías actuales con menos coordinación.
